% !PRINTED 10
\weberpage{10}

%<<CONT
und man kann darunter je zwei auswählen, bei denen das Intervall $z = 0$ bis $z = 1$ einem gegebenen der Intervalle $(x_1x_2)$, $(x_2x_3)$, $(x_3x_4)$, $(x_4x_1)$ entspricht; dem wachsenden $z$ entsprechen bei der einen dieser Transformationen die wachsenden $x$, bei der anderen die abnehmenden $x$ (mit dem Durchgange von $+\infty$ zu $-\infty$). Im ersten Falle haben die Quadratwurzeln in (9) beiderseits das gleiche, im zweiten das entgegengesetzte Zeichen.

Wenn man nicht darauf besteht, daß $\varkappa^2$ ein positiver echter Bruch ist, so kann man die Punkte $x_1$, $x_2$, $x_3$, $x_4$ auf alle Arten permutieren und erhält 24 verschiedene lineare Transformationen in die Normalform, von denen je vier dasselbe $\varkappa^2$ ergeben; man erhält im ganzen sechs verschiedene Modulen, die nach § 2, (6) auseinander abgeleitet werden.

Man übersieht am leichtesten die Gesamtheit dieser Transformationen, wenn man annimmt, das zu transformierende Integral habe bereits die Normalform:

\[
\dfrac{dx}{\sqrt{x(1-x)(1-\lambda^2x)}};
\]

man hat dann in den Formeln (6) bis (9) die $x_1$, $x_2$, $x_3$, $x_4$ auf alle möglichen Arten durch $0$, $1$, $1:\lambda^2$, $\infty$ zu ersetzen. Wir nehmen als Beispiel die folgenden Zuordnungen:

\[
\begin{array}{ccccc}
 & x_1, & x_2, & x_3, & x_4, \\\hline
1) & 0, & \infty, & \dfrac{1}{\lambda^2}, & 1, \\[10pt]
2) & 0, & \dfrac{1}{\lambda^2}, & 1, & \infty, \\[10pt]
3) & 1, & 0, & \infty, & \dfrac{1}{\lambda^2}.
\end{array}
\]

Man erhält dann aus (6) und (8)

\[
\begin{array}{c@{\quad}ll}
1) & z = \dfrac{-x}{1-x}, & \varkappa^2 = 1-\lambda^2 = \lambda'^2, \\[10pt]
2) & z = \lambda^2x, & \varkappa^2 = \dfrac{1}{\lambda^2}, \\[10pt]
3) & z = \dfrac{1-x}{1-\lambda^2x}, & \varkappa^2 = \lambda^2
\end{array}
\]

und für das Differential

\[
\dfrac{dz}{\sqrt{z(1-z)(1-\varkappa^2z)}}
\]
%CONT>>
